\documentclass[main.tex]{subfiles}
\begin{document}

\section{Price Coherence}
\label{sec:oa-price-coherence}

This appendix explains why merchants almost never surcharge despite legal ability to do so.
It covers surcharging law, empirical evidence on differential pricing, and a theoretical framework for price coherence.

I extend the baseline model to allow merchants to set different prices for card and cash consumers.
In this extension, I show that surcharges cannot steer consumers from card to cash.
Card-preferring consumers will not switch to cash unless the surcharge exceeds the merchant fee.
Second, since surcharges do not shift payment behavior, the envelope theorem makes the profit loss from uniform pricing second-order in fee dispersion.
Third, returns to surcharging across card networks are even smaller, since fee differences across networks are an order of magnitude below the cash-card gap.

Cash discounts have been legal since 1981 yet remain rare, providing a revealed-preference bound.
Menu costs suffice to deter surcharging even at the cash-card margin.
Since card-vs-card fee gaps are far smaller, those same menu costs overwhelm any incentive to price differently across networks.

\subsection{Legal Background}
\label{subsec:oa-surcharging-history}

Cash discounts have long been legal in the U.S., but card surcharges became permitted only gradually.\footnote{Under full information, discounts and surcharges are the same. But if shrouded, cash discounts are giveaways while surcharges are add-on prices \parencite{Bourguignon.etal2019}.}
The Cash Discount Act of 1981 guarantees merchants' right to offer cash discounts \parencite{Chakravorti.Shah2001, Levitin2005, Prager.etal2009}.
The Durbin Amendment in 2010 granted merchants the right to offer debit discounts \parencite{Schuh.etal2011, Briglevics.Shy2014}.
A 2013 settlement between Visa, MC, and the DOJ removed network-level no-surcharge rules, allowing surcharging in 40 states \parencite{Blakeley.Fagan2015}.\footnote{Visa allowed multi-state merchants to surcharge in states that permitted it \parencite{Visa2013}.
AmEx and Discover relaxed their no-surcharge rules to match \parencite{Merchant2016}.}
% CONSTANT: cited legal/historical fact from the 2013 settlement period
As of 2023, only Massachusetts and Connecticut ban surcharging \parencite{CardX2023}, though disclosure rules in New York and Maine make it impractical.\footnote{In New York and Maine, retailers must disclose both the credit card price and cash price in dollar amounts.
In New York, this rule ensures consumers ``should not have to do math to figure out whether they are paying the surcharge.'' \parencite{Westchester2019}.}

\subsection{Empirical Evidence}
\label{subsec:oa-surcharging-dcpc}

Surcharging is rare regardless of legality, confirming that menu costs rather than legal bans explain price coherence.
I restrict the sample to cash, check, debit, and credit transactions at retail merchants.
In the full DCPC sample (all transaction types, before this restriction), about $5\%$ of transactions involve card surcharges or cash discounts, supporting the price coherence assumption.
% HARDCODED[source: unknown, from DCPC surcharge/discount computation] current=5

I flag violations of price coherence using DCPC variables on discounts, surcharges, and payment steering.
The figure below plots three time series.
``Discount'' is the share of cash/check transactions with a discount.
``Surcharge'' is the share of credit card transactions with a fee.
``Steered'' are cash transactions by credit-preferring consumers who switched due to a discount, as a share of credit transactions plus such switched transactions.
States are grouped by surcharging legality: ``Legal'' (never banned), ``Illegal'' (still banned as of 2020), and ``Grey Area'' (repealed bans during 2013--2020).

\includegraphics[width=0.8\linewidth]{../output/graphs/price_coherence_ts.pdf}

Rates are low across all groups, and legality does not explain the rarity.
Cash discounts, always legal, are still rare.
Surcharge rates do not vary with state legal status, ruling out legal bans as the main driver of price coherence.
This pattern shows that menu costs deter differential pricing even at the cash-card margin, where fee gaps are largest.

\subsection{Private Incentives to Surcharge}
\label{subsec:oa-surcharge-theory}

\subsubsection{No Steering Between Cash and Card}

I show merchants cannot steer consumers from card to cash in a natural extension of the baseline model that permits surcharging.
The key assumption is that card convenience benefits are constant across purchases.
A consumer who values cards at merchant type $\gamma$ gets the same proportional benefit regardless of the item, ruling out within-trip substitution on payment-specific prices.


Consumers choose effective consumption $q\left(\omega\right)$ of each variety, a linear aggregate of card consumption $c\left(\omega\right)$ and cash consumption $a\left(\omega\right)$.
Merchants charge card consumers a price $1+s\left(\omega\right)$ higher.
Consumers solve
\begin{align}
U & =\max_{c,a}\left(\int_{0}^{1}q\left(\omega\right)^{\frac{\sigma-1}{\sigma}}\d\omega\right)^{\frac{\sigma}{\sigma-1}}\label{eq:ces-microfoundation-pay-choice}\\
\text{s.t. }q\left(\omega\right) & =\left(1+\gamma\left(\omega\right)v_{M\left(\omega\right)}^{w}\right)^{\frac{1}{\sigma-1}}c\left(\omega\right)+a\left(\omega\right)\\
y & \geq\int_{0}^{1}\left(c\left(\omega\right)\left(1+s\left(\omega\right)\right)+a\left(\omega\right)\right)p\left(\omega\right) \d\omega
\end{align}
Linear aggregation makes card goods higher quality and perfect substitutes for cash goods.
The baseline model corresponds to corner solutions where consumers use only their preferred payment type.

\begin{lem}
At a merchant of type $\gamma$ that accepts cards, a card consumer will use cash only if $s>\left(1+\gamma\right)^{\frac{1}{\sigma-1}}-1$\label{lem:required-surcharge}

\end{lem}
\begin{proof}
The first-order conditions for card spending $c$ and cash spending $a$ yield

\begin{align*}
\frac{\partial\mathcal{L}}{\partial c} & =I^{\frac{1}{\sigma-1}}\times q^{-\frac{1}{\sigma}}\times\left(1+\gamma\right)^{\frac{1}{\sigma-1}}-\lambda\left(1+s\right)p\\
\frac{\partial\mathcal{L}}{\partial a} & =I^{\frac{1}{\sigma-1}}\times q^{-\frac{1}{\sigma}}-\lambda p
\end{align*}
Both payment methods are used if $\left(1+\gamma\right)^{\frac{1}{\sigma-1}}=1+s$.
Since the aggregator is linear, consumers use cash exclusively if $s>\left(1+\gamma\right)^{\frac{1}{\sigma-1}}-1$.

\end{proof}

\begin{theorem}
In a market with a monopoly credit card network that charges a merchant fee of $\tau$, no merchant that accepts the credit card in the baseline model can steer consumers by setting $s=\tau$.

\end{theorem}
\begin{proof}
The lowest type accepting credit cards satisfies $\gamma^{*}=\frac{\sigma\tau}{1-\sigma\tau}$.
For general $\gamma>0,\sigma>1$, the inequality $\left(1+\gamma\right)^{\frac{1}{\sigma-1}}\geq1+\frac{\gamma}{\gamma+1}\frac{1}{\sigma-1}$ holds.
By Lemma \ref{lem:required-surcharge}, the required surcharge to induce switching exceeds $s^{*}\geq\tau\frac{\sigma}{\sigma-1}>\tau$, so setting $s=\tau$ does not induce steering.

\end{proof}

The analysis assumes consumers use only their primary card type: credit users do not switch to debit, and vice versa.
This is consistent with empirical evidence \parencite{Conrath2014} and antitrust analysis \parencite{Jones2001}.

\subsubsection{Second-Order Losses from Uniform Pricing}

Since surcharges do not change the payment method, uniform pricing generates only second-order losses by the envelope theorem.
I quantify these with a second-order Taylor approximation.

Suppose merchants charge wallet-specific prices $p^{w}$ (stacked into a vector).
After normalization, total profits are proportional to

\begin{align*}
\widehat{\Pi} & \propto\sum_{w\in\mathcal{W}}\mu^{w}\pi^{w}\\
\pi^{w} & =\left(1+\gamma v_{M}^{w}\right)\left(p^{w}\right)^{-\sigma}\left(p^{w}\left(1-\tilde{\tau}^{w}\right)-1\right)
\end{align*}
where $\tilde{\tau}^{w}\equiv\frac{(1+\gamma)\tau_M^w}{1+\gamma v_M^w}$ is the wallet-specific effective fee derived from the merchant problem (Section~\ref{subsec:appendix-quasiprofits}). The notation distinguishes it from the raw average $\tau_M^w$.
Let $p^{*}$ denote optimal payment-specific prices and $\hat{p}$ denote the optimal uniform price.

\begin{theorem}
The percentage loss from uniform pricing is:

\[
\log\widehat{\Pi}\left(p^{*}\right)-\log\widehat{\Pi}\left(\hat{p}\right)=\sum_{w}\frac{\mu^{w}\left(1+\gamma v_{M}^{w}\right)}{\sum_{l}\mu^{l}\left(1+\gamma v_{M}^{l}\right)}\times\frac{\sigma\left(\sigma-1\right)}{2}\left(\tilde{\tau}^{w}-\hat{\tilde{\tau}}\right)^{2}+O\left(\tau^{3}\right)
\]
\end{theorem}

\begin{proof}
The percentage loss is a weighted sum across wallets.
By the envelope theorem, the first-order term vanishes.
Computing the second derivative of log profit with respect to $\log p$ and evaluating at the optimal price $p^{w*}=\frac{\sigma}{\sigma-1}\left(1-\tilde{\tau}^{w}\right)^{-1}$ yields $\frac{\partial^{2}\log\pi^{w}}{\partial\left(\log p\right)^{2}}=-\sigma\left(\sigma-1\right)$.
Since $\log p^{w*}-\log\hat{p}\approx\tilde{\tau}^{w}-\hat{\tilde{\tau}}+O(\tau^2)$, the second-order Taylor expansion yields the result.

\end{proof}

High fees alone do not make uniform pricing costly. Fee \emph{dispersion} among accepted payment methods does so.
At the estimated $\sigma=$\scalarinput{param_est_merchant_ces_pass}, losses from uniform pricing are \scalarinput{surcharge_loss_multiprod} bp of profit.

\subsubsection{Card-vs-Card Surcharging}
\label{subsubsec:oa-surcharge-card-vs-card}

Merchants also seek the right to surcharge cards differently, e.g., charge AmEx more than Visa \parencite{Conrath2014}.
The results above bound the gains: in a type-symmetric equilibrium, fee dispersion across networks of the same type is near zero, so losses from uniform pricing are negligible.

The rarity of cash discounts sharpens this point.
Cash discounts have been legal since 1981, yet merchants almost never offer them.
The cash-card fee gap is roughly 2--3\% (the full merchant discount fee, since cash carries no merchant fee).
% CONSTANT: approximate interchange fee range, from industry data
That merchants do not surcharge even at this margin shows that menu costs exceed the gains at a 2--3\% fee gap. These costs include maintaining payment-specific prices, communicating them, and absorbing the resulting customer friction \parencite{Caddy.etal2020}.

Fee differences across card networks are an order of magnitude smaller.
Even away from exact type-symmetry, AmEx-Visa fee gaps are tens of bp, well below the cash-card gap.
Since merchants do not price differently at the larger margin, they will not do so at the smaller one.

\subsubsection{Limitations}
\label{subsubsec:surcharge-effective}

The results above assume every card transaction generates incremental sales for the merchant.
If some card categories do not boost sales, merchants would face stronger incentives to surcharge those payment methods.
In the limit where no card transaction generates incremental sales, surcharging at the full merchant fee would be profitable, and declining the card entirely would yield equivalent profits.

\end{document}